\scanpage{94}\mspage{11}
% archive foliation: 93
Mais avant, un commentaire. Quand on
donne $\underline{J}$, rev. cubique étale de $S$, et une
section $t$ de $\Sigma_{\underline{J}}^*$, alors la donnée d'une
courbe elliptique relative sur $S$, correspondant à
ces données $J,t$, équivaut à celle d'un
module inv. $T$ sur $S$, et d'un isom \unsure{(a)|(5)}
\[
\unsure{(a)|(5)}\qquad T^{\otimes2}\simeq T_{\Sigma,t}
\]
Quand on a disposé déjà d'un tel $T$, soit
$T_0$, alors la catégorie des $T$ \inserted{sur $S$} munis de \unsure{(a)|(5)} équivaut
à celle des $\unsure{\boldsymbol{\mu}|\mu}_{2S}$-torseurs $\underline{\varepsilon}$ sur $S$, moyennant
l'équivalence
\[
(55)\qquad \underline{\varepsilon}\longmapsto T=T_0\wedge_{\unsure{\boldsymbol{\mu}|\mu}_{2S}}\underline{\varepsilon}
\]
Ici, on a envie de prendre
\[
(56)\qquad T_0=\underline{O}_{\Sigma}(1)(t)^{\unsure{\sharp}}
=\underline{O}_{\Sigma}(1)(t)\wedge_{\unsure{\boldsymbol{\mu}|\mu}_4}Q_{0S}
\]
mais on a un iso. can.
\[
T_0^{\otimes2}\simeq\underline{O}_{\Sigma}(2)(t)\wedge_{\unsure{\boldsymbol{\mu}|\mu}_2}\unsure{\boldsymbol{\mu}|\mu}_4^*
\simeq T_{\Sigma,t}\wedge_{\unsure{\boldsymbol{\mu}|\mu}_2}(\unsure{\boldsymbol{\mu}|\mu}_4^*\wedge\underline{\omega}_{\underline{J}})
\]
et pour avoir un iso. ``raisonnable'' $T_0^{\otimes2}\simeq T_{\Sigma,t}$,
il faut se donner un \deleted{section} iso
\[
\alpha\in\unsure{\mathcal{I}}_{\unsure{\boldsymbol{\mu}|\mu}_4^*\simeq\underline{\omega}_{\underline J}}
\quad\text{i.e. une section de }\unsure{\boldsymbol{\mu}|\mu}_4^*\wedge\underline{\omega}_{\underline J}
\]
(ce qui est une des données indiquées par un
épinglage de Legendre d'une courbe elliptique $/S$
correspondant à $\underline{J}\ \ldots$). Une fois $\alpha$ fixé, on
